\itemhead{1.2 思考 · 未编号题组（教材p.14）}{XJ-B1-C01-TH-P014-Y1066}
\label{XJ-B1-C01-TH-P014-Y1066}
{\footnotesize 湘教版教材p.14 / PDF第21页 / 1.2 思考 / 未编号题组}
\field{原题}{在初中，我们已经对命题有了初步认识。你能说出命题的定义并列举出一些数学中的命题吗？}
\field{参考答案}{命题是可以判断真假的陈述句。例如“2是偶数”为真命题，“1是偶数”为假命题；未限定 $x$ 的“$x>0$”不是命题。}
{\small \field{关联知识}{L01 命题、真假、逆命题与否定（教材帮p.30 / PDF39）}
\field{方法与易错点}{明确论域与断言；真命题给理由，假命题给满足题设而使结论失败的反例。 易错：把逆命题当成否定；只改变结论中的个别词而改变原逻辑范围。}
\field{答案依据}{答案册未提供本栏目的独立对应答案；依据原题、教材概念补写参考回应。开放性总结不设唯一数值答案。}}
